\documentclass[a4paper,12pt]{article}

\usepackage[italian]{babel}
\usepackage{amsmath}
\usepackage{amsfonts}
\usepackage{graphicx}
\usepackage{listings}
\usepackage{xcolor}

\title{Implementazione e Analisi dell'Algoritmo KNN in Python}
\author{}
\date{}

\lstset{
	language=Python,
	basicstyle=\ttfamily\small,
	keywordstyle=\color{blue},
	commentstyle=\color{gray},
	showstringspaces=false
}

\begin{document}
	
	\maketitle
	
	\section{Introduzione}
	
	Il programma implementa un classificatore K-Nearest Neighbors (KNN) per prevedere la qualità del vino a partire da caratteristiche chimiche.
	
	Le fasi principali sono:
	\begin{enumerate}
		\item Caricamento dati
		\item Suddivisione train/test
		\item Normalizzazione
		\item Calcolo distanze
		\item Classificazione
		\item Valutazione
	\end{enumerate}
	
	---
	
	\section{Caricamento dei dati}
	
	\begin{lstlisting}
		def carica_dati(percorso):
		df = pd.read_csv(percorso, sep=";")
		X = df.drop(columns="quality").values
		y = df["quality"].values
		return X, y
	\end{lstlisting}
	
	Separa:
	\[
	X = \text{feature}, \quad y = \text{target}
	\]
	
	---
	
	\section{Normalizzazione Min-Max}
	
	\begin{lstlisting}
		def normalizza(X_train, X_test):
		X_min = X_train.min(axis=0)
		X_max = X_train.max(axis=0)
		X_train_norm = (X_train - X_min) / (X_max - X_min)
		X_test_norm  = (X_test  - X_min) / (X_max - X_min)
		return X_train_norm, X_test_norm
	\end{lstlisting}
	
	\subsection{Formula}
	
	\[
	x_{norm} = \frac{x - x_{min}}{x_{max} - x_{min}}
	\]
	
	\subsection{Esempio}
	
	\[
	X_{train} =
	\begin{bmatrix}
		4 & 10 \\
		6 & 20 \\
		8 & 30
	\end{bmatrix}
	\]
	
	\[
	X_{train}^{norm} =
	\begin{bmatrix}
		0 & 0 \\
		0.5 & 0.5 \\
		1 & 1
	\end{bmatrix}
	\]
	
	---
	
	\section{Distanza euclidea}
	
	\begin{lstlisting}
		def distanza_euclidea(a, b):
		return np.sqrt(np.sum((a - b) ** 2))
	\end{lstlisting}
	
	\[
	d(a,b) = \sqrt{\sum (a_i - b_i)^2}
	\]
	
	\subsection{Esempio}
	
	\[
	a = [0,0], \quad b = [1,1]
	\]
	
	\[
	d = \sqrt{2}
	\]
	
	---
	
	\section{Algoritmo KNN}
	
	\begin{lstlisting}
		def knn_predict(X_train, y_train, x_new, k):
		distanze = [distanza_euclidea(x_new, x) for x in X_train]
		indici_vicini = np.argsort(distanze)[:k]
		etichette_vicini = y_train[indici_vicini]
		valori, conteggi = np.unique(etichette_vicini, return_counts=True)
		return valori[np.argmax(conteggi)]
	\end{lstlisting}
	
	---
	
	\subsection{Esempio completo}
	
	\[
	X_{train} =
	\begin{bmatrix}
		1 & 1 \\
		2 & 2 \\
		3 & 3 \\
		6 & 6
	\end{bmatrix}
	\quad
	y_{train} = [0,0,1,1]
	\]
	
	\[
	x_{new} = [2.5,2.5], \quad k=3
	\]
	
	\subsubsection{Distanze}
	
	\[
	d = [2.12, 0.71, 0.71, 4.95]
	\]
	
	\subsubsection{Ordinamento}
	
	\[
	\text{argsort} = [1,2,0,3]
	\]
	
	\subsubsection{Vicini}
	
	\[
	[1,2,0]
	\]
	
	\subsubsection{Etichette}
	
	\[
	[0,1,0]
	\]
	
	\subsubsection{Conteggi}
	
	\[
	valori = [0,1], \quad conteggi = [2,1]
	\]
	
	\subsubsection{Risultato}
	
	\[
	\boxed{0}
	\]
	
	---
	
	\section{Valutazione del modello}
	
	\begin{lstlisting}
		previsioni = [knn_predict(X_train, y_train, x, k) for x in X_test]
		corretti = sum(p == v for p, v in zip(previsioni, y_test))
		return corretti / len(y_test)
	\end{lstlisting}
	
	---
	
	\subsection{Spiegazione}
	
	\begin{itemize}
		\item Genera tutte le previsioni
		\item Confronta con i valori reali
		\item Conta quelle corrette
		\item Calcola l’accuratezza
	\end{itemize}
	
	---
	
	\subsection{Esempio}
	
	\[
	y_{test} = [1,0,1,1,0]
	\]
	
	\[
	previsioni = [1,1,1,0,0]
	\]
	
	\[
	[True, False, True, False, True]
	\]
	
	\[
	corretti = 3
	\]
	
	\[
	\text{accuratezza} = \frac{3}{5} = 0.6
	\]
	
	---
	
	\section{Scelta del parametro K}
	
	\[
	K = [1,3,5,7,9]
	\]
	
	\begin{itemize}
		\item K piccolo → overfitting
		\item K grande → underfitting
	\end{itemize}
	
	---
	
	\section{Conclusione}
	
	Il programma implementa correttamente un modello KNN completo:
	
	\begin{itemize}
		\item preprocessing corretto
		\item distanza euclidea
		\item classificazione per maggioranza
		\item valutazione quantitativa
	\end{itemize}
	
	La normalizzazione è il passaggio più critico per garantire correttezza geometrica del modello.
	
\end{document}